This section introduces the proposed Modality Decoupling and Coupling Network (MDCN) and details its key components. As shown in Figure \ref{fig:2}, a pretrained CLIP \cite{CLIP_2021} is used as the backbone for text feature extraction. MDCN first introduces Context Prompt Decoupling Learning (CPDL), which semantically decouples modality-contextual text prompts into person-related and modality-related prompts and uses modality orthogonality to eliminate cross-modal feature redundancy. This effective decoupling ensures that modality information is clearly distinguished from identity information. Then, MDCN distinguishes modality-specific, modality-shared, and modality-contextual features through three network branches constrained by Semantic Coupling Learning (SCL). High-level semantic information is embedded via modality coupling and modality-contextual alignment and effectively guide the image encoder to focus on key attributes. Meanwhile, common ReID losses are used to boost the discriminative power of the learned features. Finally, MDCN introduces Modality Relational Knowledge Distillation (MRKD), which uses hierarchical cosine angular distillation to distill valuable high-level semantic relationships from modality-specific and modality-shared branches into modality-contextual image features. This enhances the modality-aware distinguishability in modality-contextual representation and ensures a more refined representation of cross-modal information.

Formally, we define the VI-ReID dataset as 
$\mathcal{D}_I = \{\mathcal{X}_v, \mathcal{X}_r, \mathcal{Y}_v, \mathcal{Y}_r\}$,  
where $\mathcal{X}_v = \{x_v^i\}_{i=1}^{N_v}$ and $\mathcal{X}_r = \{x_r^i\}_{i=1}^{N_r}$ denote the sets of visible and infrared images, respectively. Here, $x_v^i$ and $x_r^i$ represent the $i$-th visible and infrared samples, while $N_v$ and $N_r$ indicate the corresponding sample counts of the two modalities. The identity label sets are given as $\mathcal{Y}_v = \{y_v^i\}_{i=1}^{N_v}$ and $\mathcal{Y}_r = \{y_r^i\}_{i=1}^{N_r}$, which share a unified identity class space across modalities. 



This section introduces the proposed Modality Decoupling and Coupling Network (MDCN) and details its key components. As shown in Figure \ref{fig:2}, a pretrained CLIP \cite{CLIP_2021} is used as the backbone for text feature extraction. MDCN first introduces Context Prompt Decoupling Learning (CPDL), which semantically decouples modality-contextual text prompts into person-related and modality-related prompts and uses modality orthogonality to eliminate cross-modal feature redundancy. This effective decoupling ensures that modality information is clearly distinguished from identity information. Then, MDCN distinguishes modality-specific, modality-shared, and modality-contextual features through three network branches constrained by Semantic Coupling Learning (SCL). High-level semantic information is embedded via modality coupling and modality-contextual alignment and effectively guide the image encoder to focus on key attributes. Meanwhile, common ReID losses are used to boost the discriminative power of the learned features. Finally, MDCN introduces Modality Relational Knowledge Distillation (MRKD), which uses hierarchical cosine angular distillation to distill valuable high-level semantic relationships from modality-specific and modality-shared branches into modality-contextual image features. This enhances the modality-aware distinguishability in modality-contextual representation and ensures a more refined representation of cross-modal information.

Formally, we define the VI-ReID dataset as 
$\mathcal{D}_I = \{\mathcal{X}_v, \mathcal{X}_r, \mathcal{Y}_v, \mathcal{Y}_r\}$,  
where $\mathcal{X}_v = \{x_v^i\}_{i=1}^{N_v}$ and $\mathcal{X}_r = \{x_r^i\}_{i=1}^{N_r}$ denote the sets of visible and infrared images, respectively. Here, $x_v^i$ and $x_r^i$ represent the $i$-th visible and infrared samples, while $N_v$ and $N_r$ indicate the corresponding sample counts of the two modalities. The identity label sets are given as $\mathcal{Y}_v = \{y_v^i\}_{i=1}^{N_v}$ and $\mathcal{Y}_r = \{y_r^i\}_{i=1}^{N_r}$, which share a unified identity class space across modalities. 